\documentclass[10pt,a4paper]{amsart}
\usepackage{amsmath,amssymb,amsthm,graphicx}
\usepackage{verbatim}
\newenvironment{reference}{\comment}{\endcomment}
%\newenvironment{reference}{}{}
\newenvironment{slogan}{\comment}{\endcomment}
\newenvironment{history}{\comment}{\endcomment}

% For commutative diagrams we use Xy-pic
\usepackage[all]{xy}

% We use 2cell for 2-commutative diagrams.
\xyoption{2cell}
% Theorem environments.
%
\theoremstyle{plain}
\newtheorem{theorem}[subsection]{Theorem}
\newtheorem{proposition}[subsection]{Proposition}
\newtheorem{lemma}[subsection]{Lemma}

\theoremstyle{definition}
\newtheorem{definition}[subsection]{Definition}
\newtheorem{example}[subsection]{Example}
\newtheorem{exercise}[subsection]{Exercise}
\newtheorem{situation}[subsection]{Situation}

\theoremstyle{remark}
\newtheorem{remark}[subsection]{Remark}
\newtheorem{remarks}[subsection]{Remarks}

\numberwithin{equation}{subsection}

% Macros
%
\def\lim{\mathop{\mathrm{lim}}\nolimits}
\def\colim{\mathop{\mathrm{colim}}\nolimits}
\def\Spec{\mathop{\mathrm{Spec}}}
\def\Hom{\mathop{\mathrm{Hom}}\nolimits}
\def\Ext{\mathop{\mathrm{Ext}}\nolimits}
\def\SheafHom{\mathop{\mathcal{H}\!\mathit{om}}\nolimits}
\def\SheafExt{\mathop{\mathcal{E}\!\mathit{xt}}\nolimits}
\def\Sch{\mathit{Sch}}
\def\Mor{\mathop{\mathrm{Mor}}\nolimits}
\def\Ob{\mathop{\mathrm{Ob}}\nolimits}
\def\Sh{\mathop{\mathit{Sh}}\nolimits}
\def\NL{\mathop{N\!L}\nolimits}
\def\CH{\mathop{\mathrm{CH}}\nolimits}
\def\proetale{{pro\text{-}\acute{e}tale}}
\def\etale{{\acute{e}tale}}
\def\QCoh{\mathit{QCoh}}
\def\Ker{\mathop{\mathrm{Ker}}}
\def\Im{\mathop{\mathrm{Im}}}
\def\Coker{\mathop{\mathrm{Coker}}}
\def\Coim{\mathop{\mathrm{Coim}}}

% Boxtimes
%
\DeclareMathSymbol{\boxtimes}{\mathbin}{AMSa}{"02}

%
% Macros for moduli stacks/spaces
%
\def\QCohstack{\mathcal{QC}\!\mathit{oh}}
\def\Cohstack{\mathcal{C}\!\mathit{oh}}
\def\Spacesstack{\mathcal{S}\!\mathit{paces}}
\def\Quotfunctor{\mathrm{Quot}}
\def\Hilbfunctor{\mathrm{Hilb}}
\def\Curvesstack{\mathcal{C}\!\mathit{urves}}
\def\Polarizedstack{\mathcal{P}\!\mathit{olarized}}
\def\Complexesstack{\mathcal{C}\!\mathit{omplexes}}
% \Pic is the operator that assigns to X its picard group, usage \Pic(X)
% \Picardstack_{X/B} denotes the Picard stack of X over B
% \Picardfunctor_{X/B} denotes the Picard functor of X over B
\def\Pic{\mathop{\mathrm{Pic}}\nolimits}
\def\Picardstack{\mathcal{P}\!\mathit{ic}}
\def\Picardfunctor{\mathrm{Pic}}
\def\Deformationcategory{\mathcal{D}\!\mathit{ef}}
\usepackage{cleveref}
\crefname{originalreference}{original reference}{original references}
\Crefname{originalreference}{Original reference}{Original references}
\makeatletter
\newlabel{lemma-genus-goes-down}{{0CE3}{source}{}{sourcepointer-1}{}}
\newlabel{lemma-genus-goes-down@cref}{{[originalreference][0][]0CE3}{[1][0][]source}{}{}{}}
\makeatother
\crefname{originallemma}{original lemma}{original lemmas}
\Crefname{originallemma}{Original lemma}{Original lemmas}
\makeatletter
\newlabel{varieties-lemma-chi-tensor-finite}{{0AYT}{source}{}{sourcepointer-2}{}}
\newlabel{varieties-lemma-chi-tensor-finite@cref}{{[originallemma][0][]0AYT}{[1][0][]source}{}{}{}}
\makeatother
\crefname{originallemma}{original lemma}{original lemmas}
\Crefname{originallemma}{Original lemma}{Original lemmas}
\makeatletter
\newlabel{varieties-lemma-finite-in-codim-1}{{0AB7}{source}{}{sourcepointer-3}{}}
\newlabel{varieties-lemma-finite-in-codim-1@cref}{{[originallemma][0][]0AB7}{[1][0][]source}{}{}{}}
\makeatother
\crefname{originallemma}{original lemma}{original lemmas}
\Crefname{originallemma}{Original lemma}{Original lemmas}
\makeatletter
\newlabel{varieties-lemma-modification-normal-iso-over-codimension-1}{{0BFP}{source}{}{sourcepointer-4}{}}
\newlabel{varieties-lemma-modification-normal-iso-over-codimension-1@cref}{{[originallemma][0][]0BFP}{[1][0][]source}{}{}{}}
\makeatother
\title{Genus drops under birational proper morphisms, with a rigidity statement}
\author{The Stacks Project Authors}
\date{stacks-project:curves.tex:local-snapshot-20261012; GNU FDL 1.2 or later}
\begin{document}
\maketitle
\noindent\textit{Editorial note.} All delivered mathematical text is the author's own, copied from the pinned Stacks Project chapter file sources/stacks-curves/pinned-original-source.tex; no statement or proof step was rewritten, reordered or removed.\par\medskip
\noindent\textit{Editorial note.} The counted proof is the author's complete argument: \par\medskip
\noindent\textit{Editorial note.} The label lemma-genus-normalization names the counted statement, and the proof environment immediately following it in the source is delivered verbatim as its proof.\par\medskip
\noindent\textit{Editorial note.} The counted statement is carried by the author's own lemma environment in the chapter source.\par\medskip
\noindent\textit{Allowed standard prerequisites (editorial).}\begin{itemize}
\item Proper birational morphisms are finite in codimension one. Varieties, Lemmas \textbackslash{}label\{varieties-lemma-finite-in-codim-1\} and \textbackslash{}label\{varieties-lemma-modification-normal-iso-over-codimension-1\}. Proof lines 3889-3893: a proper birational morphism f : X' -> X is finite and an isomorphism over a dense open.
\item Genus decreases under finite birational maps. Same chapter, \textbackslash{}label\{lemma-genus-goes-down\}. Proof lines 3893-3894: the inequality dim\_k H\textasciicircum{}1(X, O\_X) >= dim\_k H\textasciicircum{}1(X', O\_\{X'\}) follows from this lemma.
\item Sheaves with zero-dimensional support are generated by global sections. Varieties, Lemma \textbackslash{}label\{varieties-lemma-chi-tensor-finite\}. Proof lines 3901-3904: the cokernel F of O\_X -> f\_*O\_\{X'\} has zero-dimensional support, hence is generated by global sections, so H\textasciicircum{}0(X, F) = 0 forces F = 0.
\end{itemize}\medskip
\ifdefined\hypertarget\hypertarget{sourcepointer-1}{}\fi
\noindent\textit{Source reference.} Original reference 0CE3: sources/stacks-curves/pinned-original-source.tex, source lines 3850--3850. Prerequisite pointer invoked at proof line(s) 3894; it names a standard fact of the same chapter that the counted argument uses, and every step of the counted proof built on it is retained verbatim in the delivered body.\par\medskip
\ifdefined\hypertarget\hypertarget{sourcepointer-2}{}\fi
\noindent\textit{Source reference.} Original lemma 0AYT: sources/stacks-chow/tags.txt, source lines 13801--13801. Prerequisite pointer invoked at proof line(s) 3904; it names a standard fact of a sibling chapter that the counted argument uses, and every step of the counted proof built on it is retained verbatim in the delivered body.\par\medskip
\ifdefined\hypertarget\hypertarget{sourcepointer-3}{}\fi
\noindent\textit{Source reference.} Original lemma 0AB7: sources/stacks-chow/tags.txt, source lines 13010--13010. Prerequisite pointer invoked at proof line(s) 3892; it names a standard fact of a sibling chapter that the counted argument uses, and every step of the counted proof built on it is retained verbatim in the delivered body.\par\medskip
\ifdefined\hypertarget\hypertarget{sourcepointer-4}{}\fi
\noindent\textit{Source reference.} Original lemma 0BFP: sources/stacks-chow/tags.txt, source lines 14392--14392. Prerequisite pointer invoked at proof line(s) 3893; it names a standard fact of a sibling chapter that the counted argument uses, and every step of the counted proof built on it is retained verbatim in the delivered body.\par\medskip
\begin{lemma}
\label{lemma-genus-normalization}
Let $k$ be a field. Let $X$ be a proper scheme of dimension $\leq 1$ over $k$.
If $X' \to X$ is a birational proper morphism, then
$$
\dim_k H^1(X, \mathcal{O}_X) \geq \dim_k H^1(X', \mathcal{O}_{X'})
$$
If $X$ is reduced, $H^0(X, \mathcal{O}_X) \to H^0(X', \mathcal{O}_{X'})$
is surjective, and equality holds, then $X' = X$.
\end{lemma}

\begin{proof}
If $f : X' \to X$ is proper birational, then it is finite and
an isomorphism over a dense open (see
Varieties, Lemmas \ref{varieties-lemma-finite-in-codim-1} and
\ref{varieties-lemma-modification-normal-iso-over-codimension-1}).
Thus the inequality by Lemma \ref{lemma-genus-goes-down}.
Assume $X$ is reduced. Then $\mathcal{O}_X \to f_*\mathcal{O}_{X'}$
is injective and we obtain a short exact sequence
$$
0 \to \mathcal{O}_X \to f_*\mathcal{O}_{X'} \to \mathcal{F} \to 0
$$
Under the assumptions given in the second statement,
we conclude from the long exact cohomology sequence that
$H^0(X, \mathcal{F}) = 0$. Then $\mathcal{F} = 0$ because
$\mathcal{F}$ is generated by global sections
(Varieties, Lemma \ref{varieties-lemma-chi-tensor-finite}).
and $\mathcal{O}_X = f_*\mathcal{O}_{X'}$. Since $f$ is affine
this implies $X = X'$.
\end{proof}

\end{document}
